\newpage
\section*{ТЕХНІЧНЕ ЗАВДАННЯ}
\addcontentsline{toc}{section}{ТЕХНІЧНЕ ЗАВДАННЯ}

\subsection*{1. Найменування та область використання}

Найменування розробки: система діалогової взаємодії з великими мовними моделями з web- та консольним інтерфейсами. Система може використовуватися як навчальний програмний продукт для демонстрації принципів побудови LLM-застосунків, роботи з локальними мовними моделями, організації persistence-рівня, web-інтерфейсу та автоматизованого тестування.

\subsection*{2. Підстава для розробки}

Підставою для розробки є завдання на виконання курсової роботи з дисципліни «Інженерія програмного забезпечення» для спеціальності 123 «Комп'ютерна інженерія».

\subsection*{3. Мета та призначення розробки}

Метою роботи є розробка програмної системи, яка забезпечує користувачу можливість вести діалог з локальною великою мовною моделлю, зберігати історію чатів, керувати моделями під час виконання програми та перевіряти роботу системи автоматизованими тестами.

Призначення системи:
\begin{itemize}
  \item створення, перегляд, перейменування та видалення чатів;
  \item надсилання повідомлень і отримання відповідей від локальної LLM;
  \item збереження історії діалогів у PostgreSQL;
  \item робота через web-інтерфейс і CLI;
  \item потокове відображення відповіді моделі;
  \item зупинка активної генерації без очікування повного завершення відповіді;
  \item пошук повідомлень і експорт діалогу у Markdown;
  \item відображення Markdown, блоків коду та математичних формул;
  \item перемикання локальних моделей і використання LoRA/QLoRA-адаптерів;
  \item вибір preset-ів генерації та перегляд стану активної моделі;
  \item тестування ключових сценаріїв роботи системи.
\end{itemize}

\subsection*{4. Джерела розробки}

Джерелами розробки є документація Python, Flask, PostgreSQL, Alembic, Hugging Face Transformers, PEFT, pytest, KaTeX, а також матеріали з архітектури програмного забезпечення, UML-моделювання та тестування програмних систем.

\subsection*{5. Технічні вимоги}

\textbf{Вимоги до програмного продукту:}
\begin{itemize}
  \item система повинна мати web-інтерфейс для роботи з чатами;
  \item система повинна підтримувати CLI-режим;
  \item система повинна зберігати історію чатів і повідомлень;
  \item система повинна підтримувати локальний LLM-backend;
  \item система повинна дозволяти перемикати локальні моделі під час роботи;
  \item система повинна підтримувати streaming generation;
  \item система повинна дозволяти зупиняти активну генерацію;
  \item система повинна підтримувати пошук повідомлень у чатах;
  \item система повинна підтримувати експорт діалогу у Markdown;
  \item система повинна відображати Markdown, блоки коду та математичні формули;
  \item система повинна показувати активний backend, модель, adapter і preset генерації;
  \item система повинна підтримувати generation presets;
  \item система повинна підтримувати підключення LoRA/QLoRA-адаптерів через конфігураційний файл;
  \item система повинна контрольовано повідомляти про помилки завантаження моделі або adapter-а;
  \item система повинна мати автоматизовані тести для основної бізнес-логіки.
\end{itemize}

\textbf{Вимоги до програмного забезпечення:}
\begin{itemize}
  \item Python 3.10 або вище;
  \item PostgreSQL;
  \item Flask;
  \item Hugging Face Transformers;
  \item PEFT для підключення LoRA/QLoRA-адаптерів;
  \item Alembic;
  \item pytest і coverage.py;
  \item KaTeX для відображення математичних формул у web-інтерфейсі.
\end{itemize}

\textbf{Вимоги до апаратної частини:}
\begin{itemize}
  \item персональний комп'ютер або ноутбук, здатний запускати Python-застосунок;
  \item достатній обсяг оперативної пам'яті для запуску обраної локальної LLM;
  \item достатній обсяг дискового простору для збереження локальних моделей і adapter-ів;
  \item GPU є бажаним для швидшої генерації, але система може працювати і на CPU з малими моделями.
\end{itemize}

\subsection*{6. Етапи розробки}

Основними етапами розробки є отримання теми, вивчення джерел, формування вимог, проєктування архітектури та UML-діаграм, реалізація програмного продукту, тестування, оформлення пояснювальної записки та підготовка до захисту.
